Table \ref{tab:classification-result} presented Top-3 and Top-1 metrics obtained from the classification outcomes, which are categorized into four sections, namely baseline, geometric transformation, color-space augmentation, and cutout-related augmentation. The training process was carried out over 200 epochs and implemented a learning rate of 0.0001. The highest metric in each group is emphasized in bold. Following the overall results, we proceed to analyze the outcomes of each group of augmentations, namely Geometric Augmentation, Color-Space Augmentation, Cutout, and RandAugment.

\begin{table}[!h]
\caption{Result of augmentation on EfficinetNet B0 at fine-tuned learning rate $10^{-5}$}    \label{tab:classification-result}

    \centering
    \small
    \tabcolsep=0.11cm
    \begin{tabular}{|l|lll|lll|}
    \hline
    & \multicolumn{3}{c}{\textbf{Top-3 Metrics}} & \multicolumn{3}{c}{\textbf{Top-1 Metrics}} \\ \hline
        \textbf{Method} & \textbf{\makecell{F1-score \\Top-3}} & \textbf{\makecell{Accuracy \\ Top-3}} & \textbf{\makecell{Precision \\ Top-3}} & \textbf{\makecell{F1-score \\ Top-1}} & \textbf{\makecell{Accuracy \\ Top-1}} & \textbf{\makecell{Precision \\ Top-1}} \\ \hline
        Baseline & 0.3434 & 0.4977 & 0.3231 & 0.2784 & 0.3339 & 0.2264 \\ \hline
        Rotation & \textbf{0.4423} & \textbf{0.5876} & \textbf{0.432} & \textbf{0.3364} & \textbf{0.4586} & \textbf{0.2585} \\ 
        Flipping & 0.4208 & 0.5826 & 0.4104 & 0.3296 & 0.4351 & 0.2563 \\ 
        Perspective & 0.3577 & 0.5209 & 0.3339 & 0.286 & 0.3414 & 0.218 \\ 
        CenterCrop 0.2 & 0.3252 & 0.5111 & 0.3028 & 0.2829 & 0.3284 & 0.2204 \\ 
        CenterCrop 0.4 & 0.316 & 0.5205 & 0.3023 & 0.2807 & 0.3333 & 0.209 \\ 
        MotionBlur & 0.3347 & 0.5209 & 0.3233 & 0.3054 & 0.3486 & 0.2353 \\ \hline
        CLAHE & 0.3027 & 0.4654 & 0.2932 & 0.2569 & 0.3093 & 0.1971 \\ 
        Sharpen & 0.3393 & 0.5095 & 0.3168 & \textbf{0.306} & 0.3525 & 0.2486 \\ 
        ChannelShuffle & 0.314 & 0.4682 & 0.2945 & 0.2646 & 0.3188 & 0.2011 \\ 
        Solarize & 0.2354 & 0.4165 & 0.2173 & 0.2251 & 0.2191 & 0.1666 \\ 
        Gaussian Noise & \textbf{0.3498} & 0.5154 & \textbf{0.3289} & 0.2868 & \textbf{0.3546} & 0.2168 \\ 
        MotionBlur & 0.3347 & \textbf{0.5209} & 0.3233 & 0.3054 & 0.3486 & \textbf{0.2353} \\ 
        % Normalize & 0.3315 & 0.5123 & 0.3119 & 0.2778 & 0.3395 & 0.2152 \\ 
        ToGray & 0.2941 & 0.4637 & 0.2787 & 0.2638 & 0.3011 & 0.211 \\ 
        ToSepia & 0.329 & 0.4708 & 0.3183 & 0.2695 & 0.3436 & 0.2111 \\ \hline
        Cutout & 0.2769 & 0.4891 & 0.2167 & 0.3387 & 0.3199 & 0.3396 \\ 
        GridDropout & 0.2827 & 0.4894 & 0.2167 & 0.333 & 0.3154 & 0.3454 \\ \hline
        RandAugment & 0.2929 & 0.5243 & 0.3789 & 0.3598 & 0.3731 & 0.2174 \\ 
    \hline
    \end{tabular}
    
\end{table}

\subsubsection{Overall}
In summary, most of the augmentation techniques did not lead to an improvement in the performance of EfficientNet. However, Rotation, Flipping, Perspective, and Gaussian Noise augmentations did show positive impacts on the model's performance. In particular, Rotation outperformed all other augmentations in terms of metrics and improved the baseline F1-Score Top-3 by almost 10\%. While Gaussian Noise was the only augmentation in the color-space group that had a positive impact on the performance, the difference was relatively small, with an increase of 0.0064 in F1-Score Top-3. Finally, regularization techniques such as cutout and dropout worsened the model's performance.

\subsubsection{Geometric Augmentations}
Figure \ref{fig:class-geo} depicts the classification performance of EfficientNet B0 when using geometric augmentations. Interestingly, contrary to the findings in the \hyperref[result:dec:geo]{detection model}and our initial hypothesis, CenterCrop augmentation does not seem to improve the model's classification performance. This observation implies that the model may require contextual information from the background to accurately predict the class. For example, in the case of a flock of birds, it may indicate a specific bird class that is commonly seen in a group, such as Graugans. Additionally, the geographical background may appear more frequently in certain bird classes, such as Brandgans, which could also assist in accurate classification.

\begin{figure}[!h]
	 \centerline{\includegraphics[width=0.9\textwidth]{src/Results/classification/class-geo.png}}
	 {\caption{Performance of Efficientnet B0 with Geometric Augmentations with fined-tuned learning rate $10^{-5}$.}\label{fig:class-geo}}
\end{figure}

Alternatively, it is possible that EfficientNet is already designed to learn features at different scales and resolutions, and by reducing the information, the model may lose some crucial details for a specific class. Therefore, CenterCrop augmentation may not be as effective in improving the model's classification performance.

However, Rotation, Flipping, and Perspective augmentations demonstrate their effectiveness in helping the model learn features in different geometric positions. With rotation the the F1-score improved substantially by 0.1. 


\subsubsection{Color-Space Augmentations}

\begin{figure}[!h]
	 \centerline{\includegraphics[width=0.9\textwidth]{src/Results/classification/class-color.png}}
	 {\caption{Performance of Efficientnet B0 with Color-space Augmentations with fined-tuned learning rate $10^{-5}$.}\label{fig:class-color}}
\end{figure}
Figure \ref{fig:class-color} displays the outcomes of the impact of color-space transformation on the performance of EfficientNets. The overall effect of this type of transformation is either negligible or insignificant with the exception of Gaussian Noise, which has a minimal positive effect. A possible explanation for this observation is that the architecture of EfficientNet is inherently designed to be resilient to color variations. The model employs a feature extractor network that learns to extract meaningful features from the input image, which is then coupled with a scaling method to ensure proper normalization of the input data. This design choice enhances the model's invariance to changes in color-space, rendering color-space augmentation less effective in improving the model's performance.


\subsubsection{Cutout Augmentations and RandAugment}


The results presented in Figure \ref{fig:class-cutout} demonstrate the impact of the cutout transformations and RandAugment on the performance of EfficientNets. It can be observed that the cutout transformation causes a decrease in the network's performance due to the loss of information. Similarly, incorporating RandAugment in the training pipeline results in a decrease in the network's performance. This could be due to the extreme nature of the augmentation techniques used in RandAugment, which can potentially slow down the convergence of the network during training.

\begin{figure}[!h]
	 \centerline{\includegraphics[width=0.9\textwidth]{src/Results/classification/f1-score-other.png}}
	 {\caption{Performance of Efficientnet B0 with Cutout, GridDropout and RandAugment with fined-tuned learning rate $10^{-5}$.}\label{fig:class-cutout}}
\end{figure}








